\par\Needspace{4\baselineskip}
\begin{center}
\hypertarget{native-037-article019}{}\pdfbookmark[2]{19.}{article-019}19.
\end{center}
Die Aufgabe, aus der wahren Anomalie die Zeit abzuleiten, lässt sich mit einer Hilfstafel, wie sie sich in mehreren astronomischen Büchern findet, viel schneller lösen. Das gilt noch mehr für die umgekehrte Aufgabe. Mit Abstand am bequemsten ist die \textsc{Barker}sche Tafel, die auch dem ausgezeichneten Werk des berühmten \textsc{Olbers} (\emph{Abhandlung über die leichteste und bequemste Methode die Bahn eines Cometen zu berechnen}, Weimar 1797.) beigefügt ist. Sie enthält für alle wahren Anomalien von 0 bis $180^0$ in Schritten von 5 Bogenminuten den Wert des Ausdrucks $75\operatorname{tang}\frac{1}{2}v+25\operatorname{tang}\frac{1}{2}v^3$ unter der Bezeichnung \emph{mittlere Bewegung}. Sucht man also die zur wahren Anomalie $v$ gehörige Zeit, so teilt man die der Tafel mit dem Argument $v$ entnommene mittlere Bewegung durch $\frac{150k}{p^{\frac{3}{2}}}$. Diese Größe heißt \emph{mittlere tägliche Bewegung}. Soll umgekehrt aus der Zeit die wahre Anomalie berechnet werden, so multipliziert man die in Tagen ausgedrückte Zeit mit $\frac{150k}{p^{\frac{3}{2}}}$. Damit erhält man die mittlere Bewegung, zu der man aus der Tafel die zugehörige Anomalie entnehmen kann. Einem negativen Wert von $v$ entsprechen offenbar dieselbe mittlere Bewegung und dieselbe Zeit, jeweils mit negativem Vorzeichen. Dieselbe Tafel ist daher für negative und positive Anomalien gleichermaßen brauchbar. Zieht man es vor, anstelle von $p$ den Perihelabstand $\frac{1}{2}p=q$ zu verwenden, so wird die mittlere tägliche Bewegung durch $\frac{k\BareRoot{2812{,}5}}{q^{\frac{3}{2}}}$ ausgedrückt. Dabei ist der konstante Faktor $k\BareRoot{2812{,}5}$ $=0{,}912\,2790\,61$ und sein Logarithmus $9{,}960\,1277\,069$. — Ist die Anomalie $v$ gefunden, so bestimmt man den Radiusvektor nach der bereits oben angegebenen Formel $r=\frac{q}{\cos\frac{1}{2}v^2}\cdot$

\par\Needspace{4\baselineskip}
\begin{center}
\hypertarget{native-037-article020}{}\pdfbookmark[2]{20.}{article-020}20.
\end{center}
Differenziert man die Gleichung $\operatorname{tang}\frac{1}{2}v+\frac{1}{3}\operatorname{tang}\frac{1}{2}v^3=2t k p^{-\frac{3}{2}}$ und behandelt dabei alle Größen $v,t,p$ als veränderlich, so erhält man
\[
\frac{\mathrm{d}v}{2\cos\frac{1}{2}v^4}=2k p^{-\frac{3}{2}}\,\mathrm{d}t-3t k p^{-\frac{5}{2}}\,\mathrm{d}p
\]
oder
\[
\mathrm{d}v=\frac{k\BareRoot{p}}{r r}\,\mathrm{d}t-\frac{3t k}{2r r\BareRoot{p}}\,\mathrm{d}p.
\]
\SourceLeaf{038}Sollen die Änderungen der Anomalie $v$ in Bogensekunden angegeben werden, so sind auch beide Terme von $\mathrm{d}v$ in dieser Einheit auszudrücken. Für $k$ ist also der in Art. 6 angegebene Wert $3548^{\prime\prime}{,}188$ zu nehmen. Führt man außerdem anstelle von $p$ die Größe $\frac{1}{2}p=q$ ein, so lautet die Formel
\[
\mathrm{d}v=\frac{k\BareRoot{2q}}{r r}\,\mathrm{d}t-\frac{3k t}{r r\BareRoot{2q}}\,\mathrm{d}q,
\]
wobei für die Konstanten die Logarithmen $\log k\BareRoot{2}=3{,}700\,5215\,724$,\\ $\log 3k\BareRoot{\frac{1}{2}}=3{,}876\,6128\,315$ zu verwenden sind.

Differenziert man ferner die Gleichung $r=\frac{p}{2\cos\frac{1}{2}v^2}$, so erhält man
\[
\frac{\mathrm{d}r}{r}=\frac{\mathrm{d}p}{p}+\operatorname{tang}\frac{1}{2}v\,\mathrm{d}v
\]
oder, wenn man $\mathrm{d}v$ durch $\mathrm{d}t$ und $\mathrm{d}p$ ausdrückt,
\[
\frac{\mathrm{d}r}{r}=\left(\frac{1}{p}-\frac{3k t\operatorname{tang}\frac{1}{2}v}{2r r\BareRoot{p}}\right)\mathrm{d}p+\frac{k\BareRoot{p}\operatorname{tang}\frac{1}{2}v}{r r}\,\mathrm{d}t.
\]
Für den Koeffizienten von $\mathrm{d}p$ erhält man, wenn man für $t$ seinen durch $v$ ausgedrückten Wert einsetzt,
\[
\begin{aligned}
&\frac{1}{p}-\frac{3p\operatorname{tang}\frac{1}{2}v^2}{4r r}-\frac{p\operatorname{tang}\frac{1}{2}v^4}{4r r}\\
&\quad=\frac{1}{r}\left(\frac{1}{2}+\frac{1}{2}\operatorname{tang}\frac{1}{2}v^2-\frac{3}{2}\sin\frac{1}{2}v^2-\frac{1}{2}\sin\frac{1}{2}v^2\operatorname{tang}\frac{1}{2}v^2\right)\\
&\quad=\frac{\cos v}{2r};
\end{aligned}
\]
Für den Koeffizienten von $\mathrm{d}t$ erhält man dagegen $=\frac{k\sin v}{r\BareRoot{p}}\cdot$ Daraus folgt $\mathrm{d}r=\frac{1}{2}\cos v\,\mathrm{d}p+\frac{k\sin v}{\BareRoot{p}}\,\mathrm{d}t$ oder, wenn man $q$ anstelle von $p$ einführt,
\[
\mathrm{d}r=\cos v\,\mathrm{d}q+\frac{k\sin v}{\BareRoot{2q}}\,\mathrm{d}t.
\]
Für die Konstante ist hier der Logarithmus $\log k\BareRoot{\frac{1}{2}}=8{,}085\,0664\,436$ zu verwenden.
